\documentclass[10pt]{article}
\usepackage[T1]{fontenc}
\usepackage[utf8]{inputenc}
\usepackage[french]{babel}
\usepackage{amsmath,amssymb,mathrsfs,mathpazo}
\usepackage{geometry}
\geometry{a4paper,margin=2cm}
\usepackage{graphicx,array,tabularx,tabularray,colortbl,multirow,multicol,booktabs,enumitem,pifont,ulem,fancybox,diagbox,cancel}
\usepackage{pstricks,pstricks-add,pst-plot,pst-tree,pst-node,pst-text,pst-func,pst-math,pst-bspline,pst-3dplot,pst-eucl,pst-solides3d}
\let\pscalculate\undefined
\usepackage{pst-bezier}
\usepackage{tikz,pgfplots,tkz-tab}
\usetikzlibrary{arrows,arrows.meta,patterns,calc,positioning,trees,decorations.markings,shapes.geometric,shapes.misc,shapes.arrows,intersections,decorations.pathmorphing,decorations.pathreplacing,backgrounds,shadows,babel}
\pgfplotsset{compat=1.18}
\usepgfplotslibrary{fillbetween}
\usepackage[np]{numprint}
\usepackage{listings,makecell,pas-tableur,twoopt}
\let\barre\undefined
\usepackage[pstricks]{tablvar}
\definecolor{darkgreen}{RGB}{0,100,0}
\providecommand{\red}{\color{red}}
\providecommand{\blue}{\color{blue}}
\providecommand{\green}{\color{green}}
\providecommand{\black}{\color{black}}
\providecommand{\white}{\color{white}}
\providecommand{\gray}{\color{gray}}
\providecommand{\cyan}{\color{cyan}}
\providecommand{\magenta}{\color{magenta}}
\providecommand{\yellow}{\color{yellow}}
\providecommand{\violet}{\color{violet}}
\providecommand{\orange}{\color{orange}}
\providecommand{\euro}{\texteuro}
\providecommand{\EUR}{\texteuro}
\providecommand{\ieme}{\textsuperscript{e}}
\providecommand{\degres}{\ensuremath{^\circ}}
\providecommand{\degree}{\ensuremath{^\circ}}
\providecommand{\vv}[1]{\overrightarrow{#1}}
\providecommand{\Coord}[1]{#1}
\providecommand{\decosix}{}
\providecommand{\starredbullet}{\ensuremath{\bullet}}
\pagestyle{empty}
\setlength{\parindent}{0pt}
\usepackage{adjustbox}\DeclareUnicodeCharacter{2236}{:}\DeclareUnicodeCharacter{2190}{\ensuremath{\leftarrow}}
\providecommand{\R}{\mathbb{R}}
\providecommand{\N}{\mathbb{N}}
\providecommand{\Z}{\mathbb{Z}}
\providecommand{\Q}{\mathbb{Q}}
\providecommand{\D}{\mathbb{D}}
\providecommand{\C}{\mathbb{C}}
\providecommand{\e}{{\mathrm{e}}}
\providecommand{\vect}[1]{\overrightarrow{#1}}
\providecommand{\vectt}[1]{\overrightarrow{\mathrm{#1}}}
\providecommand{\vvt}[1]{\overrightarrow{\mathrm{#1}}}
\providecommand{\np}[2][]{{#2~#1}}
\providecommand{\vv}[1]{{\overrightarrow{#1}}}
\renewcommand{\d}{{\mathrm{d}}}
\providecommand{\degres}{{^\circ}}
\providecommand{\degre}{{^\circ}}
\providecommand{\blue}{\color{blue}}
\providecommand{\red}{\color{red}}
\providecommand{\enskip}{\quad}
\providecommand{\llbracket}{\lbrack\!\lbrack}
\providecommand{\rrbracket}{\rbrack\!\rbrack}
\providecommand{\euro}{{\text{€}}}
\providecommand{\no}{{\text{n°}}}
\providecommand{\interval}[2]{{[#1;#2]}}
\providecommand{\intervFO}[2]{{[#1;#2[}}
\providecommand{\dotfill}{\ldots}

\providecommand{\ds}{\displaystyle}
\providecommand{\textding}[1]{#1}
\providecommand{\Oij}{\(\mathrm{O};\vec{\imath},\vec{\jmath}\)}
\providecommand{\Oijk}{\(\mathrm{O};\vec{\imath},\vec{\jmath},\vec{k}\)}
\providecommand{\Ouv}{\(\mathrm{O};\vec{u},\vec{v}\)}
\providecommand{\intervFF}[2]{\ensuremath{[#1;#2]}}
\renewcommand{\euro}{\ifmmode\text{\texteuro}\else\texteuro\fi}
\DeclareUnicodeCharacter{00B0}{\ifmmode^{\circ}\else\textdegree\fi}
\DeclareUnicodeCharacter{20AC}{\ifmmode\text{\texteuro}\else\texteuro\fi}
% Énoncés : Éducation nationale. Transcription et corrigés : APMEP, Denis Vergès et contributeurs. Découpage et classement : CoopMaths, Sébastien Lozano, Éric Elter et contributeurs. Adaptation de la présentation : Maths Explorer.
% CC BY-SA : https://coopmaths.fr/www/about/
% https://coopmaths.fr/alea/static/bac/2026/tex/bac_2026_06_sujet1_etrangers_3.tex
\begin{document}
\noindent\textit{APMEP / CoopMaths. CC BY-SA. Présentation adaptée par Maths Explorer.}\par\medskip
On considère le cube ABCDEFGH représenté ci-dessous.
	
	\begin{center}
		\begin{tikzpicture}[x=(-14:30mm),y=(18:23mm),z=(90:36mm)]
			\draw (0,0,0)node[below left]{A}--
			(1,0,0)node[below right]{B}--
			(1,1,0)node[right]{C}--
			(1,1,1)node[right]{G}--
			(0,1,1)node[above]{H}--
			(0,0,1)node[above left]{E}-- cycle
			(0,0,1) --(1,0,1)node[below right]{F}--(1,0,0)
			(1,0,1) -- (1,1,1);
			\draw[dashed] (0,1,0)node[below]{D}--(0,0,0)
			(0,1,0)--(0,1,1) (0,1,0)--(1,1,0);
			\foreach \coord/\lettre/\pos in {{(0.5,0,1)}/I/{below},{(0,0.5,1)}/J/{above left},{(0,0,0.5)}/K/{left}}
			{\draw[shift={\coord}] (2pt,2pt)--(-2pt,-2pt) (2pt,-2pt)--(-2pt,2pt) (0,0)node[\pos]{\lettre} ;}
		\end{tikzpicture}
	\end{center}
	
	On munit l'espace du repère orthonormal $\left(\mathrm{A};\vectt{AB},\vectt{AD},\vectt{AE}\right)$.
	
	Soit I le milieu du segment [EF], J le milieu de [EH] et K le milieu de [AE].
	
	\begin{enumerate}
		\item Donner les coordonnées des points I, J et K.
		
		\item \begin{enumerate}
			\item Déterminer la valeur du produit scalaire $\vectt{AI}\cdot\vectt{AJ}$.
			
			\item En déduire la mesure de l'angle $\widehat{\mathrm{IAJ}}$ arrondie au dixième de degré.
		\end{enumerate}
		
		\item \begin{enumerate}
			\item Démontrer que le vecteur $\vectt{KC}$ est un vecteur normal au plan (AIJ).
			
			\item En déduire qu'une équation cartésienne du plan (AIJ) est \quad  \[x+y-\dfrac{1}{2}z=0\]
		\end{enumerate}
		
		\item Soit L le projeté orthogonal du point C sur le plan (AIJ).
		\begin{enumerate}
			\item Déterminer les coordonnées du point L.
			
			\item En déduire que la distance du point C au plan (AIJ) est égale à $\dfrac{4}{3}$.
		\end{enumerate}
		
		\item Soit M un point de la droite (FG).
		
		On admet que les coordonnées de M sont $(1 ; m ; 1)$ avec $m$ appartenant à $\R$.
		
		\begin{enumerate}
			\item Démontrer qu'une représentation paramétrique de la droite (IM) est
			\[\begin{cases}
				x=\dfrac{1}{2} s+\dfrac{1}{2} \\
				y=m s \\
				z=1
			\end{cases} \quad \text{avec } s \in \R\]
			
			\item Peut-on affirmer que les droites (IM) et (KC) sont coplanaires quelle que soit la valeur de $m$ ?
		\end{enumerate}
	\end{enumerate}

\end{document}
